\documentclass[a4paper]{article}
\usepackage[pdftex]{graphicx}
\usepackage[utf8]{inputenc}
\usepackage{enumerate}
\usepackage{siunitx}
\sisetup{locale=DE} 
\usepackage{href-ul}
\hypersetup{
	colorlinks=true,
	linkcolor=blue,
	urlcolor=blue}
\usepackage{icomma}
\usepackage{amssymb}
\usepackage{geometry}
\geometry{a4paper, top=15mm, left=15mm, right=15mm, bottom=15mm,
	headsep=10mm, footskip=12mm}
\hyphenation{Ge-wichts-kraft}

%Source: img_0122

% Start the document
\begin{document}
%Titel
{\bf Schulaufgabe aus der Physik }\\
\begin{enumerate}[1.]

%Aufgabe 1
{\bf \item Aufgabe}

\begin{enumerate}[a)]
\item Erkläre das Entstehen von Gleitreibungskraft anhand einer Skizze.
\item Nenne zwei Maßnahmen, wie man die Gleitreibung verringern kann.
\item Erkläre den Unterschied zwischen einseitigem und zweiseitigem Hebel und gib jeweils drei Beispiele für beide Arten dazu an.
\item Was versteht man unter einem Drehmoment ?
\item Wie lautet der Energieerhaltungssatz ?
\end{enumerate}

%Aufgabe 2
{\bf \item Um die Abhängigkeit der Gleitreibungskraft von der Anpresskraft zu ermitteln, ist folgende Messreihe aufgenommen worden:}

\begin{tabular}{ccccccc}
$F_N$ in N & 0 & 0,5 & 1,5 & 2,0 & 2,5 & 3,0 \\
$F_{RG}$ in N & 0 & 0,16 & 0,47 & 0,62 & 0,78 & 0,91 
\end{tabular}

\begin{enumerate}[a)]
\item Stelle die Meßreihe graphisch dar.
\item Formuliere den gesetzmäßigen Zusammenhang, der sich ergibt.
\item Berechne die Gleitreibungszahl $\mu$ aus dem Diagramm.
\end{enumerate}


%Aufgabe 5
{\bf \item Drei gleichgroße Quader sollen über einen Holztisch gezogen werden.} Überlege, welche Möglichkeiten der Anordnung es gibt, und berechne die aufzubringende Zugkraft für die beste Möglichkeit.\\

\renewcommand{\arraystretch}{1.5}
\begin{tabular}{ccccc}
\framebox[2,2 cm]{Holz} & $F_G=\SI{5,0}{\newton}$ & $\mu_{Holz-Holz}$ &=& 0,30 \\
\framebox[2,2 cm]{Eisen} & $F_G=\SI{15}{\newton}$ & $\mu_{Eisen-Holz}$ &=& 0,40 \\
\framebox[2,2 cm]{Kunststoff} & $F_G=\SI{3,0}{\newton}$ & $\mu_{Kunststoff-Holz}$ &=& 0,35 
\end{tabular}

%Aufgabe 6
{\bf \item Auf welche Höhe wurde eine Palette mit Getränkekisten mit der Gewichtskraft von $F_G=\SI{1500}{\newton}$
gehoben, wenn an ihr eine Arbeit von $W=\SI{12}{\kilo\joule}$ verrichtet wurde?}

%Aufgabe 7
{\bf \item Ein Schrank wird in $\SI{10}{\second}$ um $\SI{2,0}{\meter}$ verschoben ($\mu = 0,20 $).} Hierfür ist eine Leistung von $P=\SI{20}{\watt}$ aufzubringen. Welche Gewichtskraft besitzt der Schrank?

\end{enumerate} 
\fbox{
	\begin{minipage}{0.5\textwidth}
		Zur Lösung bitte \href{https://www.okuyakl.de/phys/p8reibL122/ll122.pdf}{hier klicken} oder den QR-Code scannen.\\
	Weitere Arbeitsblätter gibt es unter 
	
	\href{https://www.okuyakl.de}{www.okuyakl.de}
	\end{minipage}
	\hfill
	\begin{minipage}{0.4\textwidth}
		\includegraphics[width=1.5 cm]{../../viecher/zwe03}
		\includegraphics[width=3 cm]{qr122}
		\includegraphics[width=2 cm]{../../viecher/afanticon1}
		
	\end{minipage}}
\end{document}%L Ö S U N G ---------------------------------------------
